% !TeX program = xelatex
% 星宝智能儿童早教陪伴桌——比赛汇报硬件关系图
% 推荐使用 XeLaTeX 编译；也可直接上传 Overleaf。
\documentclass[UTF8,12pt]{ctexart}

\usepackage[
  paperwidth=40.64cm,
  paperheight=22.86cm,
  margin=0cm
]{geometry}
\usepackage{tikz}
\usetikzlibrary{
  arrows.meta,
  backgrounds,
  calc,
  fit,
  positioning
}
\usepackage{xcolor}
\pagestyle{empty}
\setlength{\parindent}{0pt}

% ---------- 配色 ----------
\definecolor{Ink}{HTML}{263238}
\definecolor{Muted}{HTML}{607D8B}
\definecolor{LineGray}{HTML}{CFD8DC}
\definecolor{SoftGray}{HTML}{F5F7F8}
\definecolor{CoreBlue}{HTML}{2859D6}
\definecolor{CoreBlueDark}{HTML}{173E9F}
\definecolor{DataBlue}{HTML}{2F6FED}
\definecolor{AudioPurple}{HTML}{7A52C7}
\definecolor{VisionGreen}{HTML}{2B9A66}
\definecolor{PowerRed}{HTML}{D94B4B}
\definecolor{StructureGray}{HTML}{8B99A1}
\definecolor{WarningOrange}{HTML}{E8892F}
\definecolor{CloudCyan}{HTML}{E8F7FA}
\definecolor{TouchBlue}{HTML}{EAF1FF}
\definecolor{AudioLilac}{HTML}{F3EDFC}
\definecolor{VisionMint}{HTML}{EAF7F0}
\definecolor{PowerPink}{HTML}{FCEEEE}
\definecolor{ArmCream}{HTML}{FFF4E0}
\definecolor{ArmOrange}{HTML}{E88A21}

% ---------- TikZ 样式 ----------
\tikzset{
  module/.style={
    draw=LineGray,
    line width=0.8pt,
    rounded corners=5pt,
    fill=white,
    align=center,
    text=Ink,
    inner sep=7pt,
    minimum height=1.75cm,
    font=\fontsize{10.5}{13}\selectfont
  },
  core/.style={
    module,
    draw=CoreBlueDark,
    line width=1.4pt,
    fill=CoreBlue,
    text=white,
    minimum width=5.7cm,
    minimum height=3.15cm,
    font=\fontsize{12}{15}\selectfont
  },
  data/.style={
    -{Latex[length=3.0mm,width=2.0mm]},
    draw=DataBlue,
    line width=1.25pt
  },
  audio/.style={
    -{Latex[length=3.0mm,width=2.0mm]},
    draw=AudioPurple,
    line width=1.25pt
  },
  vision/.style={
    -{Latex[length=3.0mm,width=2.0mm]},
    draw=VisionGreen,
    line width=1.25pt,
    dashed
  },
  network/.style={
    {Latex[length=3.0mm,width=2.0mm]}-{Latex[length=3.0mm,width=2.0mm]},
    draw=DataBlue,
    line width=1.25pt
  },
  power/.style={
    -{Latex[length=2.8mm,width=1.9mm]},
    draw=PowerRed,
    line width=1.35pt
  },
  action/.style={
    -{Latex[length=3.0mm,width=2.0mm]},
    draw=ArmOrange,
    line width=1.25pt
  },
  relationlabel/.style={
    fill=white,
    text=Muted,
    inner sep=2.5pt,
    rounded corners=2pt,
    font=\fontsize{7.8}{9.6}\selectfont,
    align=center
  },
  number/.style={
    circle,
    minimum size=0.78cm,
    inner sep=0pt,
    fill=CoreBlue,
    text=white,
    font=\bfseries\fontsize{10}{10}\selectfont
  }
}

\begin{document}
\mbox{}

\begin{tikzpicture}[
  remember picture,
  overlay,
  shift={(current page.south west)},
  x=1cm,
  y=1cm
]

% ---------- 页面背景与标题 ----------
\fill[white] (0,0) rectangle (40.64,22.86);
\fill[CoreBlue] (0,20.35) rectangle (0.34,22.86);
\node[anchor=north west,text=Ink,font=\bfseries\fontsize{25}{30}\selectfont]
  at (1.15,21.95) {1. 硬件设计};
\node[anchor=north west,text=Muted,font=\fontsize{10}{13}\selectfont]
  at (1.2,20.75) {星宝智能儿童早教陪伴桌 · 多模态硬件关系与交互闭环};

% 左右分栏
\draw[LineGray,line width=0.8pt] (25.55,1.0) -- (25.55,20.25);

% ---------- 左侧：硬件关系总图 ----------
\node[anchor=north west,text=CoreBlueDark,font=\bfseries\fontsize{11}{14}\selectfont]
  at (1.15,19.85) {硬件关系总图};

% 节点
\node[module,fill=CloudCyan,minimum width=6.2cm] (cloud) at (13.05,18.0)
  {\textbf{云端智能服务}\\[-1pt]
   \footnotesize ASR · 大模型服务 · TTS};

\node[module,fill=AudioLilac,minimum width=4.8cm] (mic) at (3.85,14.95)
  {\textbf{USB 麦克风}\\[-1pt]
   \footnotesize 语音采集 · 唤醒 · 插话};

\node[module,fill=VisionMint,minimum width=4.8cm] (camera) at (3.85,10.75)
  {\textbf{USB 摄像头（扩展）}\\[-1pt]
   \footnotesize 距离检测 · 桌面感知};

\node[module,fill=SoftGray,minimum width=4.8cm,minimum height=2.15cm] (child) at (3.85,6.35)
  {\textbf{儿童与实体教具}\\[-1pt]
   \footnotesize 说话 · 触摸 · 摆放色块/卡片};

\node[core] (board) at (13.05,12.75)
  {\textbf{SC171 V3 核心计算平台}\\[3pt]
   \footnotesize 交互调度 · 本地计算\\[-1pt]
   \footnotesize 多模态融合 · 网络通信};

\node[module,fill=TouchBlue,minimum width=5.15cm,minimum height=2.4cm] (screen) at (21.7,15.1)
  {\textbf{22.5 英寸触控屏}\\[-1pt]
   \footnotesize 角色 · 内容 · 游戏\\[-1pt]
   \footnotesize 点击 · 拖拽 · 选择};

\node[module,fill=AudioLilac,minimum width=5.15cm,minimum height=2.2cm] (speaker) at (21.7,9.65)
  {\textbf{声卡 / 功放 / 扬声器}\\[-1pt]
   \footnotesize 回复 · 提示音 · 鼓励语音};

\node[module,fill=ArmCream,draw=ArmOrange,minimum width=3.85cm,minimum height=1.55cm] (armctrl) at (18.55,6.45)
  {\textbf{机械臂控制模块}\\[-1pt]
   \footnotesize 动作接收 · 映射 · 限位};

\node[module,fill=ArmCream,draw=ArmOrange,minimum width=3.75cm,minimum height=1.55cm] (arm) at (23.05,6.45)
  {\textbf{桌面机械臂}\\[-1pt]
   \footnotesize 点头 · 摇头 · 挥手 · 指向};

\node[module,fill=PowerPink,draw=PowerRed,minimum width=5.5cm,minimum height=1.45cm] (adapter)
  at (13.05,4.55)
  {\textbf{安全电源适配器 / 直流供电分配}};

% 纸箱桌体结构边界（置于背景层）
\begin{scope}[on background layer]
  \node[
    draw=StructureGray,
    line width=1.0pt,
    dashed,
    rounded corners=10pt,
    fill=SoftGray,
    fill opacity=0.27,
    inner sep=10pt,
    fit=(mic)(camera)(child)(board)(screen)(speaker)(armctrl)(arm)
  ] (deskframe) {};
\end{scope}
\node[
  anchor=south west,
  fill=white,
  text=StructureGray,
  font=\bfseries\fontsize{8.5}{10}\selectfont,
  inner sep=2pt
] at ($(deskframe.south west)+(0.35,0.10)$)
  {纸箱桌体：固定承载 · 空间分区 · 线缆整理};

% 数据与功能连接
\draw[audio] (mic.east) -- node[relationlabel,above]{USB 音频输入＋供电} (board.west |- mic.east);
\draw[vision] (camera.east) -- node[relationlabel,above]{USB 视频输入＋供电} (board.west |- camera.east);

\draw[data] ($(board.east)+(0,0.55)$) -- node[relationlabel,above]{显示信号} ($(screen.west)+(0,-0.42)$);
\draw[data] ($(screen.west)+(0,0.42)$) -- node[relationlabel,above]{触控数据回传} ($(board.east)+(0,1.15)$);

\draw[audio] (board.east |- speaker.west) -- node[relationlabel,above]{音频输出} (speaker.west);
\draw[network] (board.north) -- node[relationlabel,right]{Wi-Fi / LAN\\双向通信} (cloud.south);
\draw[action] (board.south) -- ++(0,-1.55) -| node[relationlabel,pos=0.72,above]{高层动作指令\\arm\_action} (armctrl.north);
\draw[action] (armctrl.east) -- node[relationlabel,above]{动作组执行\\安全限位} (arm.west);

% 儿童、实体桌面与输入输出的交叉关系
\draw[audio] (child.north west) to[out=125,in=-90]
  node[relationlabel,left]{儿童语音} (mic.south);
\draw[vision] (child.north east) to[out=65,in=-90]
  node[relationlabel,right]{儿童状态/教具画面} (camera.south);
\draw[data] (child.east) -- ++(1.0,0) |- node[relationlabel,pos=0.70,below]{触控操作} (screen.south);
\draw[data] (screen.south east) to[out=-45,in=20]
  node[relationlabel,right]{画面反馈} (child.east);
\draw[audio] (speaker.south west) to[out=-150,in=5]
  node[relationlabel,below]{语音反馈} (child.east);
\draw[action] (arm.south west) to[out=-165,in=5,looseness=0.85]
  node[relationlabel,below]{动作反馈 / 空间指向} (child.east);
\draw[StructureGray,dashed,line width=1.0pt]
  (arm.south) -- ++(0,-1.10)
  node[relationlabel,text=StructureGray,below]{固定在纸箱桌体};

% 电源母线及分支
\draw[PowerRed,line width=1.8pt] (6.5,3.2) -- (23.9,3.2);
\node[relationlabel,text=PowerRed,fill=white] at (8.0,3.53) {直流供电母线};
\draw[power] (adapter.south) -- (13.05,3.2);
\draw[power] (9.2,3.2) |- (board.south);
\draw[power] (20.0,3.2) |- (screen.south);
\draw[power] (22.6,3.2) |- (speaker.south);
\draw[power] (18.6,3.2) |- (armctrl.south);
\draw[power] (23.05,3.2) |- (arm.south);

% 负向耦合关系：扬声器回声进入麦克风
\draw[
  -{Latex[length=2.6mm,width=1.8mm]},
  draw=WarningOrange,
  line width=1.0pt,
  densely dotted
] (speaker.north) to[out=125,in=55,looseness=1.05]
  node[relationlabel,text=WarningOrange,above]{声学回声/串音耦合} (mic.north);

% 图例（直接绘制，避免嵌套 tikzpicture）
\draw[data] (1.15,1.15) -- (1.90,1.15);
\node[anchor=west,font=\fontsize{7.8}{9.5}\selectfont,text=Muted] at (2.02,1.15) {当前数据};
\draw[audio] (4.00,1.15) -- (4.75,1.15);
\node[anchor=west,font=\fontsize{7.8}{9.5}\selectfont,text=Muted] at (4.87,1.15) {音频};
\draw[vision] (6.15,1.15) -- (6.90,1.15);
\node[anchor=west,font=\fontsize{7.8}{9.5}\selectfont,text=Muted] at (7.02,1.15) {视觉扩展};
\draw[power] (9.35,1.15) -- (10.10,1.15);
\node[anchor=west,font=\fontsize{7.8}{9.5}\selectfont,text=Muted] at (10.22,1.15) {供电};
\draw[action] (11.85,1.15) -- (12.60,1.15);
\node[anchor=west,font=\fontsize{7.8}{9.5}\selectfont,text=Muted] at (12.72,1.15) {机械动作};
\draw[StructureGray,dashed,line width=1pt] (14.40,1.15) -- (15.15,1.15);
\node[anchor=west,font=\fontsize{7.8}{9.5}\selectfont,text=Muted] at (15.27,1.15) {结构承载};

% ---------- 右侧：核心部件及作用 ----------
\node[anchor=north west,text=CoreBlueDark,font=\bfseries\fontsize{11}{14}\selectfont]
  at (26.45,19.85) {核心部件及作用};

\node[number] at (27.0,18.05) {1};
\node[anchor=west,text width=11.5cm,align=left,text=Ink,font=\fontsize{9.7}{12.8}\selectfont]
  at (27.65,18.05)
  {\textbf{SC171 V3 核心计算平台}\\
   统一处理语音、触控、视觉与小游戏事件，并完成本地计算、交互调度和网络通信。};

\node[number] at (27.0,15.25) {2};
\node[anchor=west,text width=11.5cm,align=left,text=Ink,font=\fontsize{9.7}{12.8}\selectfont]
  at (27.65,15.25)
  {\textbf{22.5 英寸触控屏}\\
   既输出星宝形象、学习内容和游戏画面，也回传点击、拖拽、选择等触控操作。};

\node[number] at (27.0,12.45) {3};
\node[anchor=west,text width=11.5cm,align=left,text=Ink,font=\fontsize{9.7}{12.8}\selectfont]
  at (27.65,12.45)
  {\textbf{语音交互模块}\\
   USB 麦克风采集儿童语音，声卡和扬声器播放回复，形成“听取—理解—回应”闭环。};

\node[number] at (27.0,9.65) {4};
\node[anchor=west,text width=11.5cm,align=left,text=Ink,font=\fontsize{9.7}{12.8}\selectfont]
  at (27.65,9.65)
  {\textbf{视觉感知扩展模块}\\
   摄像头用于距离检测、健康提醒和后续桌面实体教具识别，不作为基础交互的强制依赖。};

\node[number] at (27.0,6.85) {5};
\node[anchor=west,text width=11.5cm,align=left,text=Ink,font=\fontsize{9.7}{12.8}\selectfont]
  at (27.65,6.85)
  {\textbf{供电与网络模块}\\
   电源系统向核心板、屏幕和音频设备分支供电；Wi-Fi/LAN 提供云端智能服务连接。};

\node[number] at (27.0,4.05) {6};
\node[anchor=west,text width=11.5cm,align=left,text=Ink,font=\fontsize{9.7}{12.8}\selectfont]
  at (27.65,4.05)
  {\textbf{纸箱桌体原型结构}\\
   完成硬件固定、线缆整理和屏幕区/教具区划分，用于验证儿童桌面交互形态。};

\node[number] at (27.0,1.75) {7};
\node[anchor=west,text width=11.5cm,align=left,text=Ink,font=\fontsize{9.2}{12.0}\selectfont]
  at (27.65,1.75)
  {\textbf{机械臂实体反馈模块}\\
   接收高层动作指令，完成挥手、点头、摇头、左右指向和小幅庆祝，与语音和屏幕形成实体陪伴反馈。};

% 页脚提示
\node[anchor=south east,text=Muted,font=\fontsize{7.2}{9}\selectfont]
  at (39.6,0.75)
  {实线：当前能力\quad 虚线：扩展能力\quad 未确认电压与接口须按实物补充};

\end{tikzpicture}
\end{document}
